\documentclass[11pt,UTF8,fontset=fandol]{ctexart}
\usepackage{amsmath,amssymb,amsfonts}
\usepackage{graphicx,float,caption,subcaption}
\usepackage{booktabs,array,multirow,makecell}
\usepackage{geometry,enumitem,xcolor}
\usepackage{listings}
\usepackage[hidelinks]{hyperref}
\usepackage{url}
\geometry{margin=2.3cm}
\graphicspath{{/Users/yye/project/论文/assets/dpo/}}
\captionsetup{font=small,labelfont=bf}
\setlength{\parskip}{0.4em}

% 数学算子
\DeclareMathOperator*{\argmax}{arg\,max}
\newcommand{\KL}{\mathrm{KL}}
\newcommand{\sref}{\mathrm{ref}}
\newcommand{\ssft}{\mathrm{SFT}}

% 代码环境样式
\lstset{
  basicstyle=\ttfamily\small,
  breaklines=true,
  columns=fullflexible,
  keepspaces=true,
  showstringspaces=false,
  frame=single,
  framesep=4pt,
  xleftmargin=4pt,
}

\title{\textbf{直接偏好优化：你的语言模型其实是一个奖励模型}\\[0.3em]
\large Direct Preference Optimization: Your Language Model is Secretly a Reward Model}
\author{%
Rafael Rafailov\textsuperscript{*\dag}\quad
Archit Sharma\textsuperscript{*\dag}\quad
Eric Mitchell\textsuperscript{*\dag}\\[0.3em]
Stefano Ermon\textsuperscript{\dag\ddag}\quad
Christopher D. Manning\textsuperscript{\dag}\quad
Chelsea Finn\textsuperscript{\dag}\\[0.4em]
\textsuperscript{\dag} Stanford University（斯坦福大学）\quad
\textsuperscript{\ddag} CZ Biohub\\[0.2em]
\texttt{\{rafailov,architsh,eric.mitchell\}@cs.stanford.edu}}
\date{第 37 届神经信息处理系统大会（NeurIPS 2023）}

\begin{document}
\maketitle

\footnotetext{\textsuperscript{*} 同等贡献；资历较浅的作者排在前面。}

\begin{abstract}
大规模无监督语言模型（LMs）虽然能够学习广博的世界知识与一定的推理能力，但由于其训练过程完全是无监督的，要实现对其行为的精确控制十分困难。现有的获得这种可操控性（steerability）的方法会收集人类对模型生成结果相对质量的标注，并对无监督语言模型进行微调，使其与这些偏好对齐，通常采用基于人类反馈的强化学习（RLHF，reinforcement learning from human feedback）。然而，RLHF 是一个复杂且往往不稳定的过程：先要拟合一个反映人类偏好的奖励模型（reward model），然后用强化学习对大型无监督语言模型进行微调，在最大化这一估计奖励的同时又不至于偏离原始模型太远。本文中，我们引入了一种对 RLHF 中奖励模型的新参数化方式，使其对应的最优策略可以以闭式（closed form）抽取出来，从而让我们能够仅用一个简单的分类损失（classification loss）来求解标准的 RLHF 问题。我们将由此得到的算法称为直接偏好优化（DPO，Direct Preference Optimization），它稳定、性能优异且计算开销轻量，免去了微调过程中从语言模型采样以及大量超参数调优的需要。我们的实验表明，DPO 微调语言模型以对齐人类偏好的效果不亚于甚至优于现有方法。值得注意的是，用 DPO 进行微调在控制生成内容情感的能力上超过了基于 PPO 的 RLHF，并在摘要与单轮对话任务中达到或提升了响应质量，同时其实现与训练都要简单得多。
\end{abstract}

\section{引言}

在超大规模数据集上训练的大型无监督语言模型（LMs）获得了令人惊讶的能力 [11, 7, 42, 8]。然而，这些模型是在由具有多种多样目标、优先级和技能水平的人类所生成的数据上训练的。其中一些目标和技能可能并不值得模仿；例如，我们也许希望我们的 AI 编程助手理解常见的编程错误，以便纠正它们，但在生成代码时，我们却希望模型偏向其训练数据中（可能稀少的）高质量编程能力。同样地，我们也许希望语言模型知道某个有 50\% 的人相信的常见误解，但我们当然不希望模型在 50\% 与之相关的查询中宣称这一误解为真！换句话说，从模型极其广博的知识与能力中选出我们所期望的响应与行为，对于构建安全、高性能且可控的 AI 系统至关重要 [28]。现有方法通常使用强化学习（RL，reinforcement learning）来引导语言模型匹配人类偏好，而我们将证明，现有方法所用的基于 RL 的目标，可以用一个简单的二元交叉熵（binary cross-entropy）目标精确地优化，从而大大简化偏好学习流程。

\begin{figure}[t]\centering
  \includegraphics[width=0.98\linewidth]{cand_p02_k00.pdf}
  \caption{\textbf{DPO 在优化人类偏好的同时避免了强化学习。}现有的利用人类反馈微调语言模型的方法，首先在一个包含提示（prompt）和人类对响应对（pairs of responses）偏好的数据集上拟合一个奖励模型，然后用 RL 寻找一个最大化所学奖励的策略。与之相反，DPO 直接用一个简单的分类目标来优化最能满足偏好的策略，拟合出一个隐式奖励模型（implicit reward model），其对应的最优策略可以以闭式抽取得到。}
  \label{fig:1}
\end{figure}

在高层次上，现有方法通过精心整理的人类偏好集合（代表人类认为安全且有帮助的行为类型）将所期望的行为灌输到语言模型中。这一偏好学习阶段发生在初始的大规模无监督预训练阶段之后，后者在一个大型文本数据集上进行。虽然偏好学习最直接的方法是在高质量响应的人类示范上进行监督微调（supervised fine-tuning），但最成功的一类方法是基于人类（或 AI）反馈的强化学习（RLHF/RLAIF；[12, 2]）。RLHF 方法在一个人类偏好数据集上拟合一个奖励模型，然后用 RL 优化语言模型策略，使其产生被赋予高奖励的响应，同时不过度偏离原始模型。虽然 RLHF 产生的模型具有令人印象深刻的对话和编程能力，但 RLHF 流程远比监督学习复杂，涉及训练多个语言模型，并在训练循环中从语言模型策略采样，带来了显著的计算成本。

本文中，我们展示了如何在不进行显式奖励建模或强化学习的情况下，直接优化语言模型以遵循人类偏好。我们提出直接偏好优化（DPO），这是一种隐式地优化与现有 RLHF 算法相同目标（带 KL 散度约束的奖励最大化）的算法，但实现简单、训练直接。直观上，DPO 的更新会提高偏好响应相对于非偏好响应的对数概率（log probability），但它引入了一个动态的、逐样本的重要性权重，以防止我们发现朴素概率比目标会导致的模型退化（model degeneration）。与现有算法一样，DPO 依赖一个理论上的偏好模型（如 Bradley-Terry 模型；[5]），用来衡量给定奖励函数与经验偏好数据的对齐程度。然而，现有方法用偏好模型定义一个偏好损失来训练奖励模型，再训练一个策略来优化所学的奖励模型；而 DPO 则通过变量替换（change of variables），将偏好损失直接定义为策略的函数。因此，给定一个关于模型响应的人类偏好数据集，DPO 可以用一个简单的二元交叉熵目标来优化策略，产生一个拟合于偏好数据的隐式奖励函数所对应的最优策略。

我们的主要贡献是直接偏好优化（DPO），一种用于从偏好中训练语言模型的简单、无 RL（RL-free）的算法。我们的实验表明，在情感调节、摘要和对话等任务中，使用参数量高达 60 亿（6B）的语言模型从偏好中学习时，DPO 至少与现有方法（包括基于 PPO 的 RLHF）一样有效。

\section{相关工作}

规模日益增大的自监督语言模型学会了零样本（zero-shot）[33] 或借助少样本（few-shot）提示 [6, 27, 11] 完成某些任务。然而，通过在指令（instructions）数据集和人类撰写的补全（completions）上微调，可以显著提升它们在下游任务上的表现以及与用户意图的对齐程度 [25, 38, 13, 41]。这种“指令微调”（instruction-tuning）过程使大语言模型（LLMs）能够泛化到指令微调集合之外的指令，并普遍提升其可用性 [13]。尽管指令微调取得了成功，但相对的人类对响应质量的判断往往比专家示范更容易收集，因此后续工作用人类偏好数据集微调 LLMs，提升了翻译 [20]、摘要 [40, 51]、故事讲述 [51] 和指令遵循 [28, 34] 方面的能力。这些方法首先在某个偏好模型（如 Bradley-Terry 模型 [5]）下优化一个神经网络奖励函数，使其与偏好数据集相容，然后用强化学习算法微调语言模型以最大化给定奖励，常用的算法有 REINFORCE [47]、近端策略优化（PPO，proximal policy optimization；[39]）或其变体 [34]。一条密切相关的研究路线利用通过人类反馈微调以遵循指令的 LLMs，为安全性或无害性（harmlessness）等特定属性生成额外的合成偏好数据 [2]，仅以文本评分标准（rubric）形式的弱监督来供 LLM 标注。这些方法代表了两个研究领域的汇聚：一个领域是用强化学习为各种目标训练语言模型 [35, 29, 48]，另一个领域是从人类偏好中学习的通用方法 [12, 21]。尽管利用相对人类偏好颇具吸引力，但用强化学习微调大型语言模型仍是一个主要的实际挑战；本工作提供了一种在理论上有依据、无需 RL 即可优化相对偏好的方法。

在语言之外，从偏好中学习策略在多臂老虎机（bandit）和强化学习两种设定下都有研究，并提出了若干方法。利用对动作的偏好或排序（而非奖励）的上下文老虎机学习，被称为上下文对决老虎机（CDB，contextual dueling bandit；[50, 14]）。在没有绝对奖励的情况下，对 CDB 的理论分析用冯·诺依曼赢家（von Neumann winner）替代了最优策略的概念，即一个对任何其他策略的期望胜率都至少为 50\% 的策略 [14]。然而，在 CDB 设定中，偏好标签是在线（online）给出的，而在从人类偏好中学习时，我们通常从一个固定批次的离线（offline）、带偏好标注的动作对中学习 [49]。类似地，基于偏好的强化学习（PbRL，preference-based RL）从由一个未知的“评分”函数（而非奖励）产生的二元偏好中学习 [9, 37]。PbRL 存在多种算法，包括能复用离策略（off-policy）偏好数据的方法，但一般都涉及先显式估计潜在的评分函数（即奖励模型），随后再对其进行优化 [16, 9, 12, 36, 21]。我们则提出一种单阶段的策略学习方法，直接优化一个策略以满足偏好。

\section{预备知识}

我们回顾 Ziegler 等人（以及后续工作 [40, 1, 28]）中的 RLHF 流程。它通常包括三个阶段：1）监督微调（SFT，supervised fine-tuning）；2）偏好采样与奖励学习；3）RL 优化。

\textbf{SFT}：RLHF 通常先用监督学习在所关注的下游任务（对话、摘要等）的高质量数据上微调一个预训练语言模型，得到一个模型 $\pi^{\ssft}$。

\textbf{奖励建模阶段}：在第二阶段，向 SFT 模型输入提示 $x$，以产生答案对 $(y_1, y_2) \sim \pi^{\ssft}(y \mid x)$。然后将这些答案对呈现给人类标注者，他们对其中一个答案表达偏好，记为 $y_w \succ y_l \mid x$，其中 $y_w$ 和 $y_l$ 分别表示 $(y_1, y_2)$ 中被偏好和不被偏好的补全。这些偏好被假定由某个潜在奖励模型 $r^*(y, x)$ 生成，而我们无法访问该模型。建模偏好有多种方法，Bradley-Terry（BT）[5] 模型是一种常用选择（如果我们能访问多个已排序的答案，更一般的 Plackett-Luce 排序模型 [32, 23] 也与该框架兼容）。BT 模型规定人类偏好分布 $p^*$ 可以写为：
\begin{equation}
p^*(y_1 \succ y_2 \mid x) = \frac{\exp\left(r^*(x, y_1)\right)}{\exp\left(r^*(x, y_1)\right) + \exp\left(r^*(x, y_2)\right)}. \tag{1}
\end{equation}

假设我们可以访问一个从 $p^*$ 采样得到的静态比较数据集 $\mathcal{D} = \left\{x^{(i)}, y_w^{(i)}, y_l^{(i)}\right\}_{i=1}^N$，我们可以参数化一个奖励模型 $r_\phi(x, y)$ 并通过极大似然（maximum likelihood）估计其参数。将该问题构造为二元分类，我们得到负对数似然损失：
\begin{equation}
\mathcal{L}_R(r_\phi, \mathcal{D}) = -\mathbb{E}_{(x, y_w, y_l) \sim \mathcal{D}}\left[\log \sigma\left(r_\phi(x, y_w) - r_\phi(x, y_l)\right)\right] \tag{2}
\end{equation}
其中 $\sigma$ 是 logistic 函数。在语言模型的语境下，网络 $r_\phi(x, y)$ 通常从 SFT 模型 $\pi^{\ssft}(y \mid x)$ 初始化，并在最后一层 transformer 层之上添加一个线性层，产生奖励值的单一标量预测 [51]。为确保奖励函数具有较低的方差，先前工作会对奖励进行归一化，使得对所有 $x$ 都有 $\mathbb{E}_{x, y \sim \mathcal{D}}[r_\phi(x, y)] = 0$。

\textbf{RL 微调阶段}：在 RL 阶段，所学的奖励函数被用来为语言模型提供反馈。遵循先前工作 [17, 18]，该优化被表述为
\begin{equation}
\max_{\pi_\theta}\ \mathbb{E}_{x \sim \mathcal{D},\, y \sim \pi_\theta(y \mid x)}\left[r_\phi(x, y)\right] - \beta\, D_{\KL}\left[\pi_\theta(y \mid x) \,\|\, \pi_\sref(y \mid x)\right], \tag{3}
\end{equation}
其中 $\beta$ 是控制偏离基准参考策略 $\pi_\sref$（即初始 SFT 模型 $\pi^{\ssft}$）程度的参数。在实践中，语言模型策略 $\pi_\theta$ 也初始化为 $\pi^{\ssft}$。所添加的约束很重要，因为它能防止模型偏离奖励模型准确的分布太远，同时维持生成的多样性并防止坍缩到单一的高奖励答案（mode-collapse）。由于语言生成的离散特性，该目标不可微，通常用强化学习来优化。标准做法 [51, 40, 1, 28] 是构造奖励函数 $r(x, y) = r_\phi(x, y) - \beta\left(\log \pi_\theta(y \mid x) - \log \pi_\sref(y \mid x)\right)$，并用 PPO [39] 进行最大化。

\section{直接偏好优化}

由于在微调语言模型这类大规模问题上应用强化学习算法存在挑战，我们的目标是推导一种直接利用偏好进行策略优化的简单方法。与先前那些先学习奖励再通过 RL 优化它的 RLHF 方法不同，我们的方法利用了一种特定的奖励模型参数化选择，使得可以以闭式抽取其最优策略，无需 RL 训练循环。正如我们接下来将详细描述的，我们的关键洞见是利用从奖励函数到最优策略的解析映射，这使我们能够把一个关于奖励函数的损失函数转化为一个关于策略的损失函数。这种变量替换方法避免了拟合一个显式、独立的奖励模型，同时仍在现有的人类偏好模型（如 Bradley-Terry 模型）下进行优化。本质上，策略网络同时表示了语言模型和（隐式的）奖励。

\textbf{推导 DPO 目标。} 我们从与先前工作相同的 RL 目标（式 3）出发，针对一般的奖励函数 $r$。遵循先前工作 [31, 30, 19, 15]，可以直接证明式 3 中带 KL 约束的奖励最大化目标的最优解具有如下形式：
\begin{equation}
\pi_r(y \mid x) = \frac{1}{Z(x)}\, \pi_\sref(y \mid x) \exp\left(\frac{1}{\beta}\, r(x, y)\right), \tag{4}
\end{equation}
其中 $Z(x) = \sum_y \pi_\sref(y \mid x) \exp\left(\frac{1}{\beta}\, r(x, y)\right)$ 是配分函数（partition function）。完整推导见附录 A.1。即使我们使用对真实奖励函数 $r^*$ 的 MLE 估计 $r_\phi$，估计配分函数 $Z(x)$ 仍然代价高昂 [19, 15]，这使得该表示在实践中难以利用。然而，我们可以重排式 4，将奖励函数表示为其对应的最优策略 $\pi_r$、参考策略 $\pi_\sref$ 和未知配分函数 $Z(\cdot)$ 的函数。具体来说，我们先对式 4 两边取对数，然后经过一些代数运算，得到：
\begin{equation}
r(x, y) = \beta \log \frac{\pi_r(y \mid x)}{\pi_\sref(y \mid x)} + \beta \log Z(x). \tag{5}
\end{equation}

我们可以把这一重参数化应用到真实奖励 $r^*$ 及其对应的最优模型 $\pi^*$ 上。幸运的是，Bradley-Terry 模型只依赖于两个补全之间的奖励之差，即 $p^*(y_1 \succ y_2 \mid x) = \sigma(r^*(x, y_1) - r^*(x, y_2))$。将式 5 中对 $r^*(x, y)$ 的重参数化代入偏好模型式 1，配分函数被消去，于是我们可以仅用最优策略 $\pi^*$ 和参考策略 $\pi_\sref$ 来表示人类偏好概率。因此，在 Bradley-Terry 模型下，最优 RLHF 策略 $\pi^*$ 满足如下偏好模型：
\begin{equation}
p^*(y_1 \succ y_2 \mid x) = \frac{1}{1 + \exp\left(\beta \log \frac{\pi^*(y_2 \mid x)}{\pi_\sref(y_2 \mid x)} - \beta \log \frac{\pi^*(y_1 \mid x)}{\pi_\sref(y_1 \mid x)}\right)} \tag{6}
\end{equation}
推导见附录 A.2。虽然式 6 使用的是 Bradley-Terry 模型，但我们也可以类似地在更一般的 Plackett-Luce 模型 [32, 23] 下推导相应表达式，见附录 A.3。

既然我们已经用最优策略（而非奖励模型）表示了人类偏好数据的概率，我们就可以为参数化策略 $\pi_\theta$ 构造一个极大似然目标。类比于奖励建模方法（即式 2），我们的策略目标变为：
\begin{equation}
\mathcal{L}_{\mathrm{DPO}}(\pi_\theta; \pi_\sref) = -\mathbb{E}_{(x, y_w, y_l) \sim \mathcal{D}}\left[\log \sigma\left(\beta \log \frac{\pi_\theta(y_w \mid x)}{\pi_\sref(y_w \mid x)} - \beta \log \frac{\pi_\theta(y_l \mid x)}{\pi_\sref(y_l \mid x)}\right)\right]. \tag{7}
\end{equation}
这样，我们用一种替代的参数化来拟合一个隐式奖励，其最优策略恰好就是 $\pi_\theta$。此外，由于我们的过程等价于拟合一个重参数化的 Bradley-Terry 模型，它享有某些理论性质，例如在偏好数据分布的适当假设下的一致性（consistencies）[4]。在第 5 节，我们将进一步讨论 DPO 与其他工作相关的理论性质。

\textbf{DPO 的更新做了什么？} 为了从机制上理解 DPO，分析损失函数 $\mathcal{L}_{\mathrm{DPO}}$ 的梯度是有益的。关于参数 $\theta$ 的梯度可以写为：
\[
\nabla_\theta \mathcal{L}_{\mathrm{DPO}}(\pi_\theta; \pi_\sref) = -\beta\, \mathbb{E}_{(x, y_w, y_l) \sim \mathcal{D}}\Big[
\underbrace{\sigma\!\left(\hat{r}_\theta(x, y_l) - \hat{r}_\theta(x, y_w)\right)}_{\text{当奖励估计有误时权重更高}}
\big[
\underbrace{\nabla_\theta \log \pi(y_w \mid x)}_{\text{提高 } y_w \text{ 的似然}}
- \underbrace{\nabla_\theta \log \pi(y_l \mid x)}_{\text{降低 } y_l \text{ 的似然}}
\big]
\Big],
\]
其中 $\hat{r}_\theta(x, y) = \beta \log \frac{\pi_\theta(y \mid x)}{\pi_\sref(y \mid x)}$ 是由语言模型 $\pi_\theta$ 和参考模型 $\pi_\sref$ 隐式定义的奖励（更多内容见第 5 节）。直观上，损失函数 $\mathcal{L}_{\mathrm{DPO}}$ 的梯度提高了被偏好补全 $y_w$ 的似然，并降低了不被偏好补全 $y_l$ 的似然。重要的是，这些样本会按照隐式奖励模型 $\hat{r}_\theta$ 给不被偏好补全打分高出多少（再由 $\beta$ 缩放）来加权，也就是隐式奖励模型对补全的排序错得有多离谱，并考虑了 KL 约束的强度。我们的实验表明这种加权的重要性，因为这种方法去掉加权系数的朴素版本会导致语言模型退化（附录表 3）。

\textbf{DPO 流程概要。} 一般的 DPO 流程如下：1）对每个提示 $x$，采样补全 $y_1, y_2 \sim \pi_\sref(\cdot \mid x)$，用人类偏好标注，构造离线偏好数据集 $\mathcal{D} = \{x^{(i)}, y_w^{(i)}, y_l^{(i)}\}_{i=1}^N$；2）针对给定的 $\pi_\sref$、$\mathcal{D}$ 和所期望的 $\beta$，优化语言模型 $\pi_\theta$ 以最小化 $\mathcal{L}_{\mathrm{DPO}}$。在实践中，人们希望复用公开可用的偏好数据集，而不是自己生成样本并收集人类偏好。由于偏好数据集是用 $\pi^{\ssft}$ 采样的，只要可用，我们就把 $\pi_\sref$ 初始化为 $\pi^{\ssft}$。然而，当 $\pi^{\ssft}$ 不可用时，我们通过最大化被偏好补全 $(x, y_w)$ 的似然来初始化 $\pi_\sref$，即 $\pi_\sref = \argmax_\pi \mathbb{E}_{x, y_w \sim \mathcal{D}}[\log \pi(y_w \mid x)]$。这一过程有助于缓解不可用的真实参考分布与 DPO 所用的 $\pi_\sref$ 之间的分布偏移（distribution shift）。关于实现和超参数的更多细节见附录 B。

\section{DPO 的理论分析}

在本节中，我们对 DPO 方法给出进一步的解释，提供理论支持，并将 DPO 的优势与用于 RLHF 的演员-评论家（actor critic）算法（如 PPO [39]）所存在的问题联系起来。

\subsection{你的语言模型其实是一个奖励模型}

DPO 能够用单一的极大似然目标，绕过拟合显式奖励和执行 RL 来学习策略这两个步骤。注意优化目标式 5 等价于一个采用奖励参数化 $r^*(x, y) = \beta \log \frac{\pi_\theta^*(y \mid x)}{\pi_\sref(y \mid x)}$ 的 Bradley-Terry 模型，而我们优化参数化模型 $\pi_\theta$，在变量替换下等价于式 2 中的奖励模型优化。在本节中，我们将构建这一重参数化背后的理论，证明它不会约束所学奖励模型的类别，并允许精确恢复最优策略。我们先定义奖励函数之间的一个等价关系。

\textbf{定义 1.} 我们称两个奖励函数 $r(x, y)$ 和 $r'(x, y)$ \emph{等价}，当且仅当对某个函数 $f$ 有 $r(x, y) - r'(x, y) = f(x)$。

容易看出这确实是一个等价关系，它将奖励函数集合划分成若干等价类。我们可以陈述以下两个引理：

\textbf{引理 1.} 在 Plackett-Luce（尤其是 Bradley-Terry）偏好框架下，来自同一类的两个奖励函数会诱导出相同的偏好分布。

\textbf{引理 2.} 来自同一等价类的两个奖励函数，在带约束的 RL 问题下会诱导出相同的最优策略。

证明很直接，我们将其推迟到附录 A.5。第一个引理是 Plackett-Luce 系列模型 [32] 中一个众所周知的欠定（under-specification）问题。由于这种欠定性，我们通常必须施加额外的可辨识性（identifiability）约束，才能对式 2 的 MLE 估计获得任何保证 [4]。第二个引理说明同一类中的所有奖励函数都产生相同的最优策略，因此对于我们最终的目标，我们只关心从最优类中恢复任意一个奖励函数。我们在附录 A.6 证明如下定理：

\textbf{定理 1.} 在温和的假设下，所有与 Plackett-Luce（尤其是 Bradley-Terry）模型相容的奖励类，都可以用重参数化 $r(x, y) = \beta \log \frac{\pi(y \mid x)}{\pi_\sref(y \mid x)}$ 来表示，其中 $\pi(y \mid x)$ 是某个模型，$\pi_\sref(y \mid x)$ 是给定的参考模型。

\textbf{证明梗概。} 考虑任意奖励函数 $r(x, y)$，它诱导出由式 4 给定的相应最优模型 $\pi_r(y \mid x)$。我们将证明 $r$ 的等价类中的一个奖励函数可以用上面给出的重参数化来表示。我们定义投影 $f$ 为
\begin{equation}
f(r; \pi_\sref, \beta)(x, y) = r(x, y) - \beta \log \sum_y \pi_\sref(y \mid x) \exp\left(\frac{1}{\beta}\, r(x, y)\right) \tag{8}
\end{equation}
算子 $f$ 只是用 $\pi_r$ 的配分函数的对数对奖励函数进行归一化。由于所添加的归一化项只是前缀 $x$ 的函数，$f(r; \pi_\sref, \beta)(x, y)$ 是 $r(x, y)$ 等价类中的一个奖励函数。最后，将 $r$ 替换为式 5 的右端（这对任意奖励函数都成立），我们得到 $f(r; \pi_\sref, \beta)(x, y) = \beta \log \frac{\pi_r(y \mid x)}{\pi_\sref(y \mid x)}$。也就是说，投影 $f$ 产生了 $r$ 等价类中具有所需形式的一个成员，并且我们提出的重参数化没有损失奖励模型的任何一般性。

我们也可以换一种方式来看待定理 1：它精确地指定了 DPO 重参数化在每个等价类中选择哪一个奖励函数，即满足如下条件的奖励函数：
\begin{equation}
\underbrace{\sum_y \pi_\sref(y \mid x) \exp\left(\frac{1}{\beta}\, r(x, y)\right)}_{=\,\pi(y \mid x),\ \text{使用定理 1 的重参数化}} = 1, \tag{9}
\end{equation}
即 $\pi(y \mid x)$ 是一个有效分布（概率非负且求和为 1）。然而，根据式 4，我们可以看出式 9 就是由奖励函数 $r(x, y)$ 诱导的最优策略的配分函数。DPO 算法的关键洞见是：我们可以对欠约束的 Plackett-Luce（尤其是 Bradley-Terry）系列偏好模型施加某些约束，从而既保留可表示奖励模型的类别，又显式地使式 4 中的最优策略对所有提示 $x$ 都在解析上可处理。

\subsection{演员-评论家算法的不稳定性}

我们也可以用我们的框架来诊断用于 RLHF 的标准演员-评论家算法（如 PPO）的不稳定性。我们遵循 RLHF 流程，聚焦于第 3 节概述的 RL 微调步骤。我们可以将第 3 节概述的带约束 RL 问题与控制即推断（control as inference）框架 [22] 联系起来。我们假设一个参数化模型 $\pi_\theta(y \mid x)$，并最小化 $D_{\KL}[\pi_\theta(y \mid x) \,\|\, \pi^*(y \mid x)]$，其中 $\pi^*$ 是由奖励函数 $r_\phi(y, x)$ 诱导的式 7 的最优策略。经过一些代数运算，这导出如下优化目标：
\begin{equation}
\max_{\pi_\theta}\ \mathbb{E}_{\pi_\theta(y \mid x)}\bigg[
r_\phi(x, y) -
\underbrace{\beta \log \sum_y \pi_\sref(y \mid x) \exp\left(\frac{1}{\beta}\, r_\phi(x, y)\right)}_{f(r_\phi,\, \pi_\sref,\, \beta)}
- \underbrace{\beta \log \frac{\pi_\theta(y \mid x)}{\pi_\sref(y \mid x)}}_{\mathrm{KL}}
\bigg] \tag{10}
\end{equation}
这与先前工作 [51, 40, 1, 28] 使用 $r_\phi$ 奖励类的 DPO 等价奖励所优化的目标相同。在此设定下，我们可以将 $f(r_\phi, \pi_\sref, \beta)$ 中的归一化项解释为参考策略 $\pi_\sref$ 的软值函数（soft value function）。虽然该项不影响最优解，但若没有它，目标的策略梯度可能具有高方差，使学习变得不稳定。我们可以用一个学习到的值函数来容纳这个归一化项，但这同样可能难以优化。另一种做法是，先前工作用人类补全基线对奖励进行归一化，这本质上是对归一化项的单样本蒙特卡洛（Monte-Carlo）估计。相比之下，DPO 的重参数化产生了一个不需要任何基线的奖励函数。

\section{实验}

在本节中，我们实证评估 DPO 直接从偏好中训练策略的能力。首先，在一个良好受控的文本生成设定中，我们提出问题：与诸如 PPO 这样的常见偏好学习算法相比，DPO 在最大化奖励与最小化相对参考策略的 KL 散度之间权衡的效率如何？接下来，我们在更大的模型和更困难的 RLHF 任务（包括摘要和对话）上评估 DPO 的表现。我们发现，在几乎不调超参数的情况下，DPO 往往表现得与诸如基于 PPO 的 RLHF 这样的强基线相当或更好，也与在学习到的奖励函数下返回 N 个采样轨迹中最优者（best of N）的做法相当。在给出这些结果之前，我们先描述实验设置；更多细节见附录 C。

\textbf{任务。} 我们的实验探索三个不同的开放式文本生成任务。在所有实验中，算法都从一个偏好数据集 $\mathcal{D} = \left\{x^{(i)}, y_w^{(i)}, y_l^{(i)}\right\}_{i=1}^N$ 中学习策略。在受控情感生成（controlled sentiment generation）中，$x$ 是来自 IMDb 数据集 [24] 的一段电影评论前缀，策略必须生成具有正面情感的 $y$。为进行受控评估，本实验中我们用一个预训练情感分类器对生成结果生成偏好对，其中 $p(\text{positive} \mid x, y_w) > p(\text{positive} \mid x, y_l)$。对于 SFT，我们在 IMDB 数据集训练集分割的评论上对 GPT-2-large 进行微调直至收敛（更多细节见附录 C.1）。在摘要（summarization）中，$x$ 是来自 Reddit 的一篇论坛帖子；策略必须生成一段概括帖子要点的摘要 $y$。遵循先前工作，我们使用 Reddit TL;DR 摘要数据集 [43] 以及 Stiennon 等人收集的人类偏好。我们使用一个在人类撰写的论坛帖子摘要上微调的 SFT 模型\footnote{\url{https://huggingface.co/CarperAI/openai_summarize_tldr_sft}}，并采用 TRLX [44] 框架进行 RLHF。人类偏好数据集是 Stiennon 等人在一个不同但训练方式类似的 SFT 模型的样本上收集的。最后，在单轮对话（single-turn dialogue）中，$x$ 是人类查询，内容可以是从天体物理学问题到情感建议请求的任何东西。策略必须对用户的查询产生一个引人入胜且有帮助的响应 $y$；我们使用 Anthropic Helpful and Harmless 对话数据集 [1]，其中包含 17 万（170k）条人类与自动助手之间的对话。每条记录以一对由一个大型（但来源未知）语言模型生成的响应结尾，并附有一个偏好标签，标明人类偏好的响应。在此设定中，没有可用的预训练 SFT 模型；因此我们仅在被偏好的补全上微调一个现成的语言模型来构成 SFT 模型。

\textbf{评估。} 我们的实验使用两种不同的评估方法。为分析每个算法在优化带约束的奖励最大化目标方面的有效性，在受控情感生成设定中，我们用每个算法所达到的奖励与相对参考策略的 KL 散度的前沿（frontier）来评估它；这一前沿之所以可计算，是因为我们可以访问真实奖励函数（一个情感分类器）。然而，在现实世界中，真实奖励函数是未知的；因此，在摘要和单轮对话设定中，我们分别用算法相对一个基线策略的胜率（win rate）来评估它们，其中以 GPT-4 作为人类评估摘要质量和响应有用性的代理（proxy）。对于摘要，我们以测试集中的参考摘要作为基线；对于对话，我们以测试数据集中被偏好的响应作为基线。虽然已有研究表明语言模型可以成为比现有指标更好的自动评估器 [10]，但我们在第 6.4 节进行了一项人类研究，以证明使用 GPT-4 进行评估的合理性。我们发现 GPT-4 的判断与人类高度相关，人类与 GPT-4 的一致程度通常与人类标注者之间的一致程度相当或更高。

\begin{figure}[t]\centering
  \includegraphics[width=0.99\linewidth]{fig2.pdf}
  \caption{\textbf{左图。}期望奖励相对于到参考策略 KL 的前沿。DPO 在所有 KL 值上都提供了最高的期望奖励，展示了其优化的质量。\textbf{右图。}TL;DR 摘要相对人类撰写摘要的胜率，以 GPT-4 作为评估器。DPO 在摘要上超过了 PPO 的最佳情况性能，同时对采样温度的变化更稳健。}
  \label{fig:2}
\end{figure}

\textbf{方法。} 除 DPO 外，我们还评估了若干现有的训练语言模型以遵循人类偏好的方法。最简单地，我们在摘要任务中探索用 GPT-J [45] 进行零样本提示，在对话任务中用 Pythia-2.8B [3] 进行 2-shot 提示。此外，我们评估 SFT 模型以及 Preferred-FT，后者是在被选中的补全 $y_w$ 上用监督学习微调的模型，$y_w$ 来自 SFT 模型（在受控情感和摘要中）或一个通用语言模型（在单轮对话中）。另一种伪监督方法是 Unlikelihood [46]，它只是优化策略以最大化赋予 $y_w$ 的概率并最小化赋予 $y_l$ 的概率；我们在“unlikelihood”项上使用一个可选系数 $\alpha \in [0, 1]$。我们还考虑了使用从偏好数据学习到的奖励函数的 PPO [39]，以及 PPO-GT，后者是一个理想化对照（oracle），它从受控情感设定中可用的真实奖励函数学习。在我们的情感实验中，我们使用 PPO-GT 的两种实现，一种是现成版本 [44]，另一种是修改版本，它对奖励进行归一化并进一步调优超参数以提升性能（在用学习到的奖励运行“常规”PPO 时我们也使用这些修改）。最后，我们考虑 Best of N 基线，从 SFT 模型（在对话中为 Preferred-FT）采样 N 个响应，并根据从偏好数据集学习到的奖励函数返回得分最高的响应。这种高性能方法将奖励模型的质量与 PPO 优化解耦，但即使对于中等大小的 N 在计算上也不切实际，因为它需要在测试时为每个查询采样 N 个补全。

\subsection{DPO 优化 RLHF 目标的效果如何？}

典型 RLHF 算法所用的带 KL 约束的奖励最大化目标，在利用奖励的同时限制策略不要偏离参考策略太远。因此，在比较算法时，我们必须同时考虑所达到的奖励以及 KL 差异；以高得多的 KL 为代价换取略高的奖励未必可取。图 2 展示了情感设定中各算法的奖励-KL 前沿。我们对每个算法执行多次训练运行，每次运行使用不同的策略保守性超参数（PPO 的目标 KL $\in \{3, 6, 9, 12\}$，$\beta \in \{0.05, 0.1, 1, 5\}$，unlikelihood 的 $\alpha \in \{0.05, 0.1, 0.5, 1\}$，Preferred-FT 用随机种子）。这一扫描总共包含 22 次运行。每训练 100 步后直至收敛，我们在一组测试提示上评估每个策略，计算真实奖励函数下的平均奖励以及相对参考策略的平均序列级 KL\footnote{即各时间步 KL 散度之和。}，即 $\KL(\pi \,\|\, \pi_\sref)$。我们发现 DPO 产生了迄今为止最高效的前沿，在仍保持低 KL 的同时达到了最高奖励。这一结果在多个方面尤为值得注意。首先，DPO 和 PPO 优化的是相同的目标，但 DPO 明显更高效；DPO 的奖励/KL 权衡严格优于 PPO。其次，即使 PPO 能访问真实奖励（PPO-GT），DPO 也比 PPO 取得了更好的前沿。

\begin{figure}[t]\centering
  \includegraphics[width=0.99\linewidth]{fig3.pdf}
  \caption{\textbf{左图。}GPT-4 为 Anthropic-HH 单步对话计算的胜率；DPO 是唯一一个在 Anthropic-HH 测试集中超越被选中摘要（chosen summaries）的方法。\textbf{右图。}训练过程中不同采样温度下的胜率。对于不同采样温度，DPO 相对数据集标签的提升在整个训练过程中相当稳定。}
  \label{fig:3}
\end{figure}

\subsection{DPO 能否扩展到真实偏好数据集？}

接下来，我们评估 DPO 在摘要和单轮对话上的微调性能。对于摘要，诸如 ROUGE 这样的自动评估指标可能与人类偏好相关性较差 [40]，并且先前工作发现用 PPO 在人类偏好上微调语言模型能提供更有效的摘要。我们通过在 TL;DR 摘要数据集的测试集分割上采样补全，并计算相对测试集中参考补全的平均胜率，来评估不同方法。所有方法的补全都在从 0.0 到 1.0 变化的温度下采样，胜率见图 2（右）。DPO、PPO 和 Preferred-FT 都微调相同的 GPT-J SFT 模型\footnote{\url{https://huggingface.co/CarperAI/openai_summarize_tldr_sft}}。我们发现 DPO 在温度 0.0 时的胜率约为 61\%，超过了 PPO 在其最优采样温度 0.0 时 57\% 的性能。DPO 还达到了比 best of N 基线更高的最大胜率。我们注意到我们没有对 DPO 的 $\beta$ 超参数做有意义的调优，因此这些结果可能低估了 DPO 的潜力。此外，我们发现 DPO 比 PPO 对采样温度更稳健，PPO 的性能在高温下可能退化到基础 GPT-J 模型的水平。Preferred-FT 相对 SFT 模型没有显著提升。我们还在第 6.4 节进行了 DPO 与 PPO 的人类正面对比评估，其中温度 0.25 的 DPO 样本相对温度 0 的 PPO 样本被偏好的比例为 58\%。

在单轮对话上，我们在 Anthropic HH 数据集 [1] 测试集分割中含一步人类-助手交互的子集上评估不同方法。GPT-4 评估以测试集中被偏好的补全作为参考来计算不同方法的胜率。由于此任务没有标准的 SFT 模型，我们从一个预训练的 Pythia-2.8B 出发，用 Preferred-FT 在被选中的补全上训练一个参考模型，使补全位于模型的分布之内，然后用 DPO 训练。我们还与 128 个 Preferred-FT 补全中的最优者进行比较（我们发现 Best of N 基线在此任务中于 128 个补全处趋于平稳；见附录图 4），以及与 Pythia-2.8B 基础模型的 2-shot 提示版本进行比较，发现对于每个方法表现最好的温度，DPO 表现得相当或更好。我们还评估了一个来自知名来源\footnote{\url{https://github.com/CarperAI/trlx/tree/main/examples/hh}}、用 PPO 在 Anthropic HH 数据集\footnote{\url{https://huggingface.co/reciprocate/ppo_hh_pythia-6B}}上训练的 RLHF 模型，但无法找到一个提示或采样温度使其性能优于基础 Pythia-2.8B 模型。基于我们在 TL;DR 上的结果以及两种方法优化相同奖励函数这一事实，我们将 Best of 128 视为 PPO 级别性能的粗略代理。总体而言，DPO 是唯一一个计算高效且能在 Anthropic HH 数据集上超越被偏好补全的方法，并提供了与计算开销巨大的 Best of 128 基线相当或更好的性能。最后，图 3 表明 DPO 相对较快地收敛到其最佳性能。

\subsection{对新输入分布的泛化}

为进一步比较 PPO 和 DPO 在分布偏移下的性能，我们将来自 Reddit TL;DR 摘要实验的 PPO 和 DPO 策略，在一个不同的分布——CNN/DailyMail 数据集 [26] 测试集分割中的新闻文章上进行评估，使用来自 TL;DR 的最佳采样温度（0 和 0.25）。结果见表 1。我们计算了相对数据集中真实摘要的 GPT-4 胜率，使用与 Reddit TL;DR 相同的 GPT-4 (C) 提示，但把其中的“forum post”（论坛帖子）一词替换为“news article”（新闻文章）。对于这一新分布，DPO 继续以显著优势超越 PPO 策略。这一实验初步证明，即使 DPO 没有使用 PPO 所用的额外无标注 Reddit TL;DR 提示，DPO 策略也能与 PPO 策略同样好地泛化。

\begin{table}[H]\centering
\caption{在分布外（out-of-distribution）的 CNN/DailyMail 输入文章上，相对真实摘要的 GPT-4 胜率。}
\label{tab:1}
\begin{tabular}{lcc}
\toprule
 & \multicolumn{2}{c}{相对真实摘要的胜率} \\
\cmidrule(lr){2-3}
算法 & Temp 0 & Temp 0.25 \\
\midrule
DPO & 0.36 & 0.31 \\
PPO & 0.26 & 0.23 \\
\bottomrule
\end{tabular}
\end{table}

\subsection{用人类判断验证 GPT-4 的判断}

我们进行了一项人类研究，以验证 GPT-4 判断的可靠性，使用了 TL;DR 摘要实验的结果和两种不同的 GPT-4 提示。GPT-4 (S)（simple，简单）提示只是询问哪份摘要更好地概括了帖子中的重要信息。GPT-4 (C)（concise，简洁）提示还会询问哪份摘要更简洁；我们评估这一提示，是因为我们发现用 GPT-4 (S) 提示时，GPT-4 偏好比人类所偏好的更长、更重复的摘要。完整提示见附录 C.2。我们进行三组比较，使用表现最高的（DPO，温度 0.25）、最低的（PPO，温度 1.0）和居中的（SFT，温度 0.25）方法，目的是覆盖多样的样本质量；三种方法都与贪婪采样的 PPO（其表现最好的温度）进行比较。我们发现，使用这两种提示时，GPT-4 与人类一致的频率大致与人类彼此一致的频率相当，表明 GPT-4 是人类评估的一个合理代理（由于人类评分者有限，我们只为 DPO 和 PPO-1 比较收集了多个人类判断）。总体而言，GPT-4 (C) 提示通常提供更能代表人类的胜率；因此我们在第 6.2 节的主要结果中使用这一提示。关于人类研究的更多细节，包括呈现给评分者的网页界面和人类志愿者名单，见附录 D.3。

\begin{table}[H]\centering
\caption{在 TL;DR 摘要样本上比较人类与 GPT-4 的胜率以及逐项判断的一致率。人类与 GPT-4 一致的程度大致与他们彼此一致的程度相当。每项实验都将所述方法的摘要与温度为 0 的 PPO 摘要进行比较。}
\label{tab:2}
\begin{tabular}{lccc}
\toprule
 & DPO & SFT & PPO-1 \\
\midrule
受访者人数 N & 272 & 122 & 199 \\
GPT-4 (S) 胜率 \% & 47 & 27 & 13 \\
GPT-4 (C) 胜率 \% & 54 & 32 & 12 \\
人类胜率 \% & 58 & 43 & 17 \\
GPT-4 (S)-H 一致率 & 70 & 77 & 86 \\
GPT-4 (C)-H 一致率 & 67 & 79 & 85 \\
H-H 一致率 & 65 & - & 87 \\
\bottomrule
\end{tabular}
\end{table}

\section{讨论}

从偏好中学习是训练有能力、对齐的语言模型的一个强大、可扩展的框架。我们引入了 DPO，一种无需强化学习、从偏好中训练语言模型的简单训练范式。DPO 不是为了使用现成的 RL 算法而把偏好学习问题强行塞入标准的 RL 设定，而是识别出语言模型策略与奖励函数之间的一种映射，使得能够用一个简单的交叉熵损失、不借助强化学习、也不损失一般性地直接训练语言模型以满足人类偏好。在几乎不调超参数的情况下，DPO 表现得与现有 RLHF 算法（包括那些基于 PPO 的算法）相当或更好；因此，DPO 切实降低了从人类偏好中训练更多语言模型的门槛。

\textbf{局限与未来工作。} 我们的结果为未来工作提出了若干重要问题。与从显式奖励函数学习相比，DPO 策略的分布外泛化（generalize out of distribution）能力如何？我们的初步结果表明 DPO 策略能与基于 PPO 的模型同样好地泛化，但需要更全面的研究。例如，用来自 DPO 策略的自标注（self-labeling）训练，是否也能有效利用无标注提示？另一方面，奖励的过度优化（reward over-optimization）在直接偏好优化设定中如何表现，图 3-右中性能的轻微下降是否就是其一个实例？此外，虽然我们评估了参数量高达 6B 的模型，但探索把 DPO 扩展到规模大几个数量级的最先进模型，是未来工作一个令人兴奋的方向。在评估方面，我们发现 GPT-4 计算的胜率受提示影响；未来工作可以研究如何最好地从自动化系统中引出高质量的判断。最后，DPO 在从人类偏好训练语言模型之外还有许多可能的应用，包括在其他模态（modalities）中训练生成模型。

\section*{致谢}
\addcontentsline{toc}{section}{致谢}

EM 感激地致谢来自 Knight-Hennessy 研究生奖学金的资助。CF 和 CM 是 CIFAR 学者。本工作部分由 Stanford Accelerator for Learning（SAL）和 Stanford Institute for Human-Centered Artificial Intelligence（HAI）的 Generative AI for the Future of Learning 种子资助项目支持。Stanford Center for Research on Foundation Models（CRFM）提供了本工作实验所用的部分计算资源。本工作部分由 ONR 资助 N00014-20-1-2675 支持。

\section*{参考文献}
\addcontentsline{toc}{section}{参考文献}

\begin{enumerate}[leftmargin=2.2em, itemsep=2pt]
\item Y. Bai, A. Jones, K. Ndousse, A. Askell, A. Chen, N. DasSarma, D. Drain, S. Fort, D. Ganguli, T. Henighan, N. Joseph, S. Kadavath, J. Kernion, T. Conerly, S. El-Showk, N. Elhage, Z. Hatfield-Dodds, D. Hernandez, T. Hume, S. Johnston, S. Kravec, L. Lovitt, N. Nanda, C. Olsson, D. Amodei, T. Brown, J. Clark, S. McCandlish, C. Olah, B. Mann, and J. Kaplan. Training a helpful and harmless assistant with reinforcement learning from human feedback, 2022.
\item Y. Bai, S. Kadavath, S. Kundu, A. Askell, J. Kernion, A. Jones, A. Chen, A. Goldie, A. Mirhoseini, C. McKinnon, C. Chen, C. Olsson, C. Olah, D. Hernandez, D. Drain, D. Ganguli, D. Li, E. Tran-Johnson, E. Perez, J. Kerr, J. Mueller, J. Ladish, J. Landau, K. Ndousse, K. Lukosuite, L. Lovitt, M. Sellitto, N. Elhage, N. Schiefer, N. Mercado, N. DasSarma, R. Lasenby, R. Larson, S. Ringer, S. Johnston, S. Kravec, S. E. Showk, S. Fort, T. Lanham, T. Telleen-Lawton, T. Conerly, T. Henighan, T. Hume, S. R. Bowman, Z. Hatfield-Dodds, B. Mann, D. Amodei, N. Joseph, S. McCandlish, T. Brown, and J. Kaplan. Constitutional ai: Harmlessness from ai feedback, 2022.
\item S. Biderman, H. Schoelkopf, Q. Anthony, H. Bradley, K. O'Brien, E. Hallahan, M. A. Khan, S. Purohit, U. S. Prashanth, E. Raff, A. Skowron, L. Sutawika, and O. van der Wal. Pythia: A suite for analyzing large language models across training and scaling, 2023.
\item H. Bong and A. Rinaldo. Generalized results for the existence and consistency of the MLE in the Bradley-Terry-Luce model. International Conference on Machine Learning, 2022. arXiv:2110.11487.
\item R. A. Bradley and M. E. Terry. Rank analysis of incomplete block designs: I. the method of paired comparisons. Biometrika, 39(3/4):324--345, 1952. doi: https://doi.org/10.2307/2334029.
\item T. Brown, B. Mann, N. Ryder, M. Subbiah, J. D. Kaplan, P. Dhariwal, A. Neelakantan, P. Shyam, G. Sastry, A. Askell, S. Agarwal, A. Herbert-Voss, G. Krueger, T. Henighan, R. Child, A. Ramesh, D. Ziegler, J. Wu, C. Winter, C. Hesse, M. Chen, E. Sigler, M. Litwin, S. Gray, B. Chess, J. Clark, C. Berner, S. McCandlish, A. Radford, I. Sutskever, and D. Amodei. Language models are few-shot learners. In H. Larochelle, M. Ranzato, R. Hadsell, M. Balcan, and H. Lin, editors, Advances in Neural Information Processing Systems, volume 33, pages 1877--1901. Curran Associates, Inc., 2020. URL \url{https://proceedings.neurips.cc/paper_files/paper/2020/file/1457c0d6bfcb4967418bfb8ac142f64a-Paper.pdf}.
\item T. Brown, B. Mann, N. Ryder, M. Subbiah, J. D. Kaplan, P. Dhariwal, A. Neelakantan, P. Shyam, G. Sastry, A. Askell, et al. Language models are few-shot learners. Advances in neural information processing systems, 33:1877--1901, 2020.
\item S. Bubeck, V. Chandrasekaran, R. Eldan, J. Gehrke, E. Horvitz, E. Kamar, P. Lee, Y. T. Lee, Y. Li, S. Lundberg, H. Nori, H. Palangi, M. T. Ribeiro, and Y. Zhang. Sparks of artificial general intelligence: Early experiments with GPT-4, 2023. arXiv preprint arXiv:2303.12712.
\item R. Busa-Fekete, B. Sz\"or\'enyi, P. Weng, W. Cheng, and E. H\"ullermeier. Preference-based reinforcement learning: evolutionary direct policy search using a preference-based racing algorithm. Machine Learning, 97(3):327--351, July 2014. doi: 10.1007/s10994-014-5458-8. URL \url{https://doi.org/10.1007/s10994-014-5458-8}.
\item Y. Chen, R. Wang, H. Jiang, S. Shi, and R.-L. Xu. Exploring the use of large language models for reference-free text quality evaluation: A preliminary empirical study. ArXiv, abs/2304.00723, 2023.
\item A. Chowdhery, S. Narang, J. Devlin, M. Bosma, G. Mishra, A. Roberts, P. Barham, H. W. Chung, C. Sutton, S. Gehrmann, et al. Palm: Scaling language modeling with pathways. arXiv preprint arXiv:2204.02311, 2022.
\item P. F. Christiano, J. Leike, T. Brown, M. Martic, S. Legg, and D. Amodei. Deep reinforcement learning from human preferences. In I. Guyon, U. V. Luxburg, S. Bengio, H. Wallach, R. Fergus, S. Vishwanathan, and R. Garnett, editors, Advances in Neural Information Processing Systems, volume 30. Curran Associates, Inc., 2017. URL \url{https://proceedings.neurips.cc/paper_files/paper/2017/file/d5e2c0adad503c91f91df240d0cd4e49-Paper.pdf}.
\item H. W. Chung, L. Hou, S. Longpre, B. Zoph, Y. Tay, W. Fedus, Y. Li, X. Wang, M. Dehghani, S. Brahma, A. Webson, S. S. Gu, Z. Dai, M. Suzgun, X. Chen, A. Chowdhery, A. Castro-Ros, M. Pellat, K. Robinson, D. Valter, S. Narang, G. Mishra, A. Yu, V. Zhao, Y. Huang, A. Dai, H. Yu, S. Petrov, E. H. Chi, J. Dean, J. Devlin, A. Roberts, D. Zhou, Q. V. Le, and J. Wei. Scaling instruction-finetuned language models, 2022.
\item M. Dud\'ik, K. Hofmann, R. E. Schapire, A. Slivkins, and M. Zoghi. Contextual dueling bandits. In P. Gr\"unwald, E. Hazan, and S. Kale, editors, Proceedings of The 28th Conference on Learning Theory, volume 40 of Proceedings of Machine Learning Research, pages 563--587, Paris, France, 03--06 Jul 2015. PMLR. URL \url{https://proceedings.mlr.press/v40/Dudik15.html}.
\item D. Go, T. Korbak, G. Kruszewski, J. Rozen, N. Ryu, and M. Dymetman. Aligning language models with preferences through f-divergence minimization. In Proceedings of the 40th International Conference on Machine Learning, ICML'23. JMLR.org, 2023.
\item A. Jain, B. Wojcik, T. Joachims, and A. Saxena. Learning trajectory preferences for manipulators via iterative improvement. In C. Burges, L. Bottou, M. Welling, Z. Ghahramani, and K. Weinberger, editors, Advances in Neural Information Processing Systems, volume 26. Curran Associates, Inc., 2013. URL \url{https://proceedings.neurips.cc/paper_files/paper/2013/file/c058f544c737782deacefa532d9add4c-Paper.pdf}.
\item N. Jaques, S. Gu, D. Bahdanau, J. M. Hern\'andez-Lobato, R. E. Turner, and D. Eck. Sequence tutor: Conservative fine-tuning of sequence generation models with kl-control. In International Conference on Machine Learning, pages 1645--1654. PMLR, 2017.
\item N. Jaques, J. H. Shen, A. Ghandeharioun, C. Ferguson, A. Lapedriza, N. Jones, S. S. Gu, and R. Picard. Human-centric dialog training via offline reinforcement learning. arXiv preprint arXiv:2010.05848, 2020.
\item T. Korbak, H. Elsahar, G. Kruszewski, and M. Dymetman. On reinforcement learning and distribution matching for fine-tuning language models with no catastrophic forgetting. In S. Koyejo, S. Mohamed, A. Agarwal, D. Belgrave, K. Cho, and A. Oh, editors, Advances in Neural Information Processing Systems, volume 35, pages 16203--16220. Curran Associates, Inc., 2022. URL \url{https://proceedings.neurips.cc/paper_files/paper/2022/file/67496dfa96afddab795530cc7c69b57a-Paper-Conference.pdf}.
\item J. Kreutzer, J. Uyheng, and S. Riezler. Reliability and learnability of human bandit feedback for sequence-to-sequence reinforcement learning. In Proceedings of the 56th Annual Meeting of the Association for Computational Linguistics (Volume 1: Long Papers), pages 1777--1788, Melbourne, Australia, July 2018. Association for Computational Linguistics. doi: 10.18653/v1/P18-1165. URL \url{https://aclanthology.org/P18-1165}.
\item A. Kupcsik, D. Hsu, and W. S. Lee. Learning Dynamic Robot-to-Human Object Handover from Human Feedback, pages 161--176. Springer International Publishing, 01 2018. ISBN 978-3-319-51531-1. doi: 10.1007/978-3-319-51532-8\_10.
\item S. Levine. Reinforcement learning and control as probabilistic inference: Tutorial and review, 2018.
\item R. D. Luce. Individual choice behavior: A theoretical analysis. Courier Corporation, 2012.
\item A. L. Maas, R. E. Daly, P. T. Pham, D. Huang, A. Y. Ng, and C. Potts. Learning word vectors for sentiment analysis. In Proceedings of the 49th Annual Meeting of the Association for Computational Linguistics: Human Language Technologies, pages 142--150, Portland, Oregon, USA, June 2011. Association for Computational Linguistics. URL \url{http://www.aclweb.org/anthology/P11-1015}.
\item S. Mishra, D. Khashabi, C. Baral, and H. Hajishirzi. Cross-task generalization via natural language crowdsourcing instructions. In Proceedings of the 60th Annual Meeting of the Association for Computational Linguistics (Volume 1: Long Papers), pages 3470--3487, Dublin, Ireland, May 2022. Association for Computational Linguistics. doi: 10.18653/v1/2022.acl-long.244. URL \url{https://aclanthology.org/2022.acl-long.244}.
\item R. Nallapati, B. Zhou, C. dos Santos, \c{C}. Gul\c{c}ehre, and B. Xiang. Abstractive text summarization using sequence-to-sequence RNNs and beyond. In Proceedings of the 20th SIGNLL Conference on Computational Natural Language Learning, pages 280--290, Berlin, Germany, Aug. 2016. Association for Computational Linguistics. doi: 10.18653/v1/K16-1028. URL \url{https://aclanthology.org/K16-1028}.
\item D. Narayanan, M. Shoeybi, J. Casper, P. LeGresley, M. Patwary, V. Korthikanti, D. Vainbrand, P. Kashinkunti, J. Bernauer, B. Catanzaro, A. Phanishayee, and M. Zaharia. Efficient large-scale language model training on gpu clusters using megatron-lm. In Proceedings of the International Conference for High Performance Computing, Networking, Storage and Analysis, SC '21, New York, NY, USA, 2021. Association for Computing Machinery. ISBN 9781450384421. doi: 10.1145/3458817.3476209. URL \url{https://doi.org/10.1145/3458817.3476209}.
\item L. Ouyang, J. Wu, X. Jiang, D. Almeida, C. Wainwright, P. Mishkin, C. Zhang, S. Agarwal, K. Slama, A. Ray, J. Schulman, J. Hilton, F. Kelton, L. Miller, M. Simens, A. Askell, P. Welinder, P. F. Christiano, J. Leike, and R. Lowe. Training language models to follow instructions with human feedback. In S. Koyejo, S. Mohamed, A. Agarwal, D. Belgrave, K. Cho, and A. Oh, editors, Advances in Neural Information Processing Systems, volume 35, pages 27730--27744. Curran Associates, Inc., 2022. URL \url{https://proceedings.neurips.cc/paper_files/paper/2022/file/b1efde53be364a73914f58805a001731-Paper-Conference.pdf}.
\item R. Paulus, C. Xiong, and R. Socher. A deep reinforced model for abstractive summarization. In International Conference on Learning Representations, 2018. URL \url{https://openreview.net/forum?id=HkAClQgA-}.
\item X. B. Peng, A. Kumar, G. Zhang, and S. Levine. Advantage-weighted regression: Simple and scalable off-policy reinforcement learning. arXiv preprint arXiv:1910.00177, 2019.
\item J. Peters and S. Schaal. Reinforcement learning by reward-weighted regression for operational space control. In Proceedings of the 24th international conference on Machine learning, pages 745--750, 2007.
\item R. L. Plackett. The analysis of permutations. Journal of the Royal Statistical Society. Series C (Applied Statistics), 24(2):193--202, 1975. doi: https://doi.org/10.2307/2346567.
\item A. Radford, J. Wu, R. Child, D. Luan, D. Amodei, and I. Sutskever. Language models are unsupervised multitask learners, 2019. Ms., OpenAI.
\item R. Ramamurthy, P. Ammanabrolu, K. Brantley, J. Hessel, R. Sifa, C. Bauckhage, H. Hajishirzi, and Y. Choi. Is reinforcement learning (not) for natural language processing: Benchmarks, baselines, and building blocks for natural language policy optimization. In The Eleventh International Conference on Learning Representations, 2023. URL \url{https://openreview.net/forum?id=8aHzds2uUyB}.
\item M. Ranzato, S. Chopra, M. Auli, and W. Zaremba. Sequence level training with recurrent neural networks. CoRR, abs/1511.06732, 2015.
\item D. Sadigh, A. D. Dragan, S. Sastry, and S. A. Seshia. Active preference-based learning of reward functions. In Robotics: Science and Systems (RSS), 2017.
\item A. Saha, A. Pacchiano, and J. Lee. Dueling rl: Reinforcement learning with trajectory preferences. In F. Ruiz, J. Dy, and J.-W. van de Meent, editors, Proceedings of The 26th International Conference on Artificial Intelligence and Statistics, volume 206 of Proceedings of Machine Learning Research, pages 6263--6289. PMLR, 25--27 Apr 2023. URL \url{https://proceedings.mlr.press/v206/saha23a.html}.
\item V. Sanh, A. Webson, C. Raffel, S. Bach, L. Sutawika, Z. Alyafeai, A. Chaffin, A. Stiegler, A. Raja, M. Dey, M. S. Bari, C. Xu, U. Thakker, S. S. Sharma, E. Szczechla, T. Kim, G. Chhablani, N. Nayak, D. Datta, J. Chang, M. T.-J. Jiang, H. Wang, M. Manica, S. Shen, Z. X. Yong, H. Pandey, R. Bawden, T. Wang, T. Neeraj, J. Rozen, A. Sharma, A. Santilli, T. Fevry, J. A. Fries, R. Teehan, T. L. Scao, S. Biderman, L. Gao, T. Wolf, and A. M. Rush. Multitask prompted training enables zero-shot task generalization. In International Conference on Learning Representations, 2022. URL \url{https://openreview.net/forum?id=9Vrb9D0WI4}.
\item J. Schulman, F. Wolski, P. Dhariwal, A. Radford, and O. Klimov. Proximal policy optimization algorithms, 2017.
\item N. Stiennon, L. Ouyang, J. Wu, D. M. Ziegler, R. Lowe, C. Voss, A. Radford, D. Amodei, and P. Christiano. Learning to summarize from human feedback, 2022.
\item R. Thoppilan, D. D. Freitas, J. Hall, N. Shazeer, A. Kulshreshtha, H.-T. Cheng, A. Jin, T. Bos, L. Baker, Y. Du, Y. Li, H. Lee, H. S. Zheng, A. Ghafouri, M. Menegali, Y. Huang, M. Krikun, D. Lepikhin, J. Qin, D. Chen, Y. Xu, Z. Chen, A. Roberts, M. Bosma, V. Zhao, Y. Zhou, C.-C. Chang, I. Krivokon, W. Rusch, M. Pickett, P. Srinivasan, L. Man, K. Meier-Hellstern, M. R. Morris, T. Doshi, R. D. Santos, T. Duke, J. Soraker, B. Zevenbergen, V. Prabhakaran, M. Diaz, B. Hutchinson, K. Olson, A. Molina, E. Hoffman-John, J. Lee, L. Aroyo, R. Rajakumar, A. Butryna, M. Lamm, V. Kuzmina, J. Fenton, A. Cohen, R. Bernstein, R. Kurzweil, B. Aguera-Arcas, C. Cui, M. Croak, E. Chi, and Q. Le. Lamda: Language models for dialog applications, 2022.
\item H. Touvron, T. Lavril, G. Izacard, X. Martinet, M.-A. Lachaux, T. Lacroix, B. Rozi\`ere, N. Goyal, E. Hambro, F. Azhar, et al. Llama: Open and efficient foundation language models. arXiv preprint arXiv:2302.13971, 2023.
\item M. V\"olske, M. Potthast, S. Syed, and B. Stein. TL;DR: Mining Reddit to learn automatic summarization. In Proceedings of the Workshop on New Frontiers in Summarization, pages 59--63, Copenhagen, Denmark, Sept. 2017. Association for Computational Linguistics. doi: 10.18653/v1/W17-4508. URL \url{https://aclanthology.org/W17-4508}.
\item L. von Werra, J. Tow, reciprocated, S. Matiana, A. Havrilla, cat state, L. Castricato, Alan, D. V. Phung, A. Thakur, A. Bukhtiyarov, aaronrmm, F. Milo, Daniel, D. King, D. Shin, E. Kim, J. Wei, M. Romero, N. Pochinkov, O. Sanseviero, R. Adithyan, S. Siu, T. Simonini, V. Blagojevic, X. Song, Z. Witten, alexandremuzio, and crumb. CarperAI/trlx: v0.6.0: LLaMa (Alpaca), Benchmark Util, T5 ILQL, Tests, Mar. 2023. URL \url{https://doi.org/10.5281/zenodo.7790115}.
\item B. Wang and A. Komatsuzaki. GPT-J-6B: A 6 Billion Parameter Autoregressive Language Model. \url{https://github.com/kingoflolz/mesh-transformer-jax}, May 2021.
\item S. Welleck, I. Kulikov, S. Roller, E. Dinan, K. Cho, and J. Weston. Neural text generation with unlikelihood training. arXiv preprint arXiv:1908.04319, 2019.
\item R. J. Williams. Simple statistical gradient-following algorithms for connectionist reinforcement learning. Mach. Learn., 8(3--4):229--256, may 1992. ISSN 0885-6125. doi: 10.1007/BF00992696. URL \url{https://doi.org/10.1007/BF00992696}.
\item Y. Wu and B. Hu. Learning to extract coherent summary via deep reinforcement learning. In Proceedings of the Thirty-Second AAAI Conference on Artificial Intelligence and Thirtieth Innovative Applications of Artificial Intelligence Conference and Eighth AAAI Symposium on Educational Advances in Artificial Intelligence, AAAI'18/IAAI'18/EAAI'18. AAAI Press, 2018. ISBN 978-1-57735-800-8.
\item X. Yan, C. Luo, C. L. A. Clarke, N. Craswell, E. M. Voorhees, and P. Castells. Human preferences as dueling bandits. In Proceedings of the 45th International ACM SIGIR Conference on Research and Development in Information Retrieval, SIGIR '22, page 567--577, New York, NY, USA, 2022. Association for Computing Machinery. ISBN 9781450387323. doi: 10.1145/3477495.3531991. URL \url{https://doi.org/10.1145/3477495.3531991}.
\item Y. Yue, J. Broder, R. Kleinberg, and T. Joachims. The k-armed dueling bandits problem. Journal of Computer and System Sciences, 78(5):1538--1556, 2012. ISSN 0022-0000. doi: https://doi.org/10.1016/j.jcss.2011.12.028. URL \url{https://www.sciencedirect.com/science/article/pii/S0022000012000281}. JCSS Special Issue: Cloud Computing 2011.
\item D. M. Ziegler, N. Stiennon, J. Wu, T. B. Brown, A. Radford, D. Amodei, P. Christiano, and G. Irving. Fine-tuning language models from human preferences, 2020.
\end{enumerate}

\section*{作者贡献}
\addcontentsline{toc}{section}{作者贡献}

所有作者都为实验的设计、分析与迭代，论文的撰写与编辑，以及项目进展的整体管理做出了宝贵贡献。

RR 在与 EM 的讨论中提出使用自回归奖励模型；推导了 DPO 目标；证明了该算法的理论性质并撰写了相关章节和附录。他还建议并协助组织了实验，并贡献了部分 PPO 和奖励学习基线。

AS 发起了关于使用加权回归方法作为 PPO 替代方案的讨论；发起了与项目相关的组织工作，撰写了将 DPO 与加权回归和 unlikelihood 联系起来的初步分析；进行了 DPO 与基线实现的设计与迭代，以及 DPO 的初步探索性实验；进行了大量的实验组织与设计（数据集、基线、评估）；主导了受控情感生成和摘要任务的模型训练与评估；进行了 GPT-4 评估的设计迭代（尤其是摘要）；为摘要、预备/方法和实验部分做出了大量撰写贡献；并为其他章节做出了编辑贡献。

EM 在关于学习自回归奖励函数的早期讨论中提供了意见；编写了 DPO 的第一个实现并运行了首批 DPO 实验；训练了论文实验中使用的大规模（摘要和对话）DPO 模型；进行了初步的 GPT-4 胜率评估并搭建了相关基础设施；为人类研究招募了参与者、组织实施并分析了结果；撰写了摘要、引言、相关工作、讨论和大部分实验内容；并协助编辑了论文的其余部分。

CF、CM 和 SE 指导了研究，提出了想法和实验，并协助撰写论文。

\appendix

\section{数学推导}

\subsection{推导带 KL 约束的奖励最大化目标的最优解}

在本附录中，我们将推导式 4。类似于式 3，我们在任意奖励函数 $r(x, y)$、参考模型 $\pi_\sref$ 和一般的非参数策略类下，优化如下目标：
\begin{equation}
\max_\pi\ \mathbb{E}_{x \sim \mathcal{D},\, y \sim \pi}\left[r(x, y)\right] - \beta\, D_{\KL}\left[\pi(y \mid x) \,\|\, \pi_\sref(y \mid x)\right] \tag{11}
\end{equation}
我们现在有：
\begin{equation}
\begin{aligned}
&\max_\pi\ \mathbb{E}_{x \sim \mathcal{D},\, y \sim \pi}\left[r(x, y)\right] - \beta\, D_{\KL}\left[\pi(y \mid x) \,\|\, \pi_\sref(y \mid x)\right] \\
&= \max_\pi\ \mathbb{E}_{x \sim \mathcal{D}}\, \mathbb{E}_{y \sim \pi(y \mid x)}\left[r(x, y) - \beta \log \frac{\pi(y \mid x)}{\pi_\sref(y \mid x)}\right] \\
&= \min_\pi\ \mathbb{E}_{x \sim \mathcal{D}}\, \mathbb{E}_{y \sim \pi(y \mid x)}\left[\log \frac{\pi(y \mid x)}{\pi_\sref(y \mid x)} - \frac{1}{\beta}\, r(x, y)\right] \\
&= \min_\pi\ \mathbb{E}_{x \sim \mathcal{D}}\, \mathbb{E}_{y \sim \pi(y \mid x)}\left[\log \frac{\pi(y \mid x)}{\frac{1}{Z(x)} \pi_\sref(y \mid x) \exp\left(\frac{1}{\beta} r(x, y)\right)} - \log Z(x)\right]
\end{aligned}
\tag{12}
\end{equation}
其中配分函数为：
\[
Z(x) = \sum_y \pi_\sref(y \mid x) \exp\left(\frac{1}{\beta}\, r(x, y)\right).
\]
注意配分函数只是 $x$ 和参考策略 $\pi_\sref$ 的函数，而不依赖于策略 $\pi$。我们现在可以定义
\[
\pi^*(y \mid x) = \frac{1}{Z(x)}\, \pi_\sref(y \mid x) \exp\left(\frac{1}{\beta}\, r(x, y)\right),
\]
它是一个有效的概率分布，因为对所有 $y$ 都有 $\pi^*(y \mid x) \geq 0$ 且 $\sum_y \pi^*(y \mid x) = 1$。由于 $Z(x)$ 不是 $y$ 的函数，我们可以将式 12 中的最终目标重新组织为：
\begin{equation}
\min_\pi\ \mathbb{E}_{x \sim \mathcal{D}}\, \mathbb{E}_{y \sim \pi(y \mid x)}\left[\log \frac{\pi(y \mid x)}{\pi^*(y \mid x)} - \log Z(x)\right] = \tag{13}
\end{equation}
\begin{equation}
\min_\pi\ \mathbb{E}_{x \sim \mathcal{D}}\left[D_{\KL}\!\left(\pi(y \mid x) \,\|\, \pi^*(y \mid x)\right) - \log Z(x)\right] \tag{14}
\end{equation}
现在，由于 $Z(x)$ 不依赖于 $\pi$，最小值由最小化第一个 KL 项的策略达到。Gibbs 不等式告诉我们，KL 散度当且仅当两个分布相同时取到最小值 0。因此我们得到最优解：
\begin{equation}
\pi(y \mid x) = \pi^*(y \mid x) = \frac{1}{Z(x)}\, \pi_\sref(y \mid x) \exp\left(\frac{1}{\beta}\, r(x, y)\right) \tag{15}
\end{equation}
对所有 $x \in \mathcal{D}$ 成立。推导至此完成。

\subsection{在 Bradley-Terry 模型下推导 DPO 目标}

在 Bradley-Terry 偏好模型下推导 DPO 目标是直接的，因为我们有
\begin{equation}
p^*(y_1 \succ y_2 \mid x) = \frac{\exp\left(r^*(x, y_1)\right)}{\exp\left(r^*(x, y_1)\right) + \exp\left(r^*(x, y_2)\right)} \tag{16}
\end{equation}
在第 4 节中我们证明了可以通过其对应的最优策略来表示（不可用的）真实奖励：
\begin{equation}
r^*(x, y) = \beta \log \frac{\pi^*(y \mid x)}{\pi_\sref(y \mid x)} + \beta \log Z(x) \tag{17}
\end{equation}
将式 17 代入式 16，我们得到：
\[
\begin{aligned}
p^*(y_1 \succ y_2 \mid x) &= \frac{\exp\left(\beta \log \frac{\pi^*(y_1 \mid x)}{\pi_\sref(y_1 \mid x)} + \beta \log Z(x)\right)}{\exp\left(\beta \log \frac{\pi^*(y_1 \mid x)}{\pi_\sref(y_1 \mid x)} + \beta \log Z(x)\right) + \exp\left(\beta \log \frac{\pi^*(y_2 \mid x)}{\pi_\sref(y_2 \mid x)} + \beta \log Z(x)\right)} \\
&= \frac{1}{1 + \exp\left(\beta \log \frac{\pi^*(y_2 \mid x)}{\pi_\sref(y_2 \mid x)} - \beta \log \frac{\pi^*(y_1 \mid x)}{\pi_\sref(y_1 \mid x)}\right)} \\
&= \sigma\left(\beta \log \frac{\pi^*(y_1 \mid x)}{\pi_\sref(y_1 \mid x)} - \beta \log \frac{\pi^*(y_2 \mid x)}{\pi_\sref(y_2 \mid x)}\right).
\end{aligned}
\]
最后一行就是式 7 中的逐实例损失。

\subsection{在 Plackett-Luce 模型下推导 DPO 目标}

Plackett-Luce 模型 [32, 23] 是 Bradley-Terry 模型在排序（而非仅成对比较）上的推广。与 Bradley-Terry 模型类似，它规定当人们面对一组可能的选择时，会以正比于某个潜在奖励函数对该选择取值的概率来偏好该选择。在我们的语境中，当面对一个提示 $x$ 和一组 $K$ 个答案 $y_1, \ldots, y_K$ 时，用户会输出一个排列 $\tau: [K] \to [K]$，给出他们对这些答案的排序。Plackett-Luce 模型规定
\begin{equation}
p^*(\tau \mid y_1, \ldots, y_K, x) = \prod_{k=1}^{K} \frac{\exp\left(r^*(x, y_{\tau(k)})\right)}{\sum_{j=k}^{K} \exp\left(r^*(x, y_{\tau(j)})\right)} \tag{18}
\end{equation}
注意当 $K = 2$ 时，式 18 退化为 Bradley-Terry 模型。然而，对于一般的 Plackett-Luce 模型，我们仍然可以利用式 5 的结果，代入由其最优策略参数化的奖励函数。类似于附录 A.2，归一化常数 $Z(x)$ 被消去，我们得到：
\begin{equation}
p^*(\tau \mid y_1, \ldots, y_K, x) = \prod_{k=1}^{K} \frac{\exp\left(\beta \log \frac{\pi^*(y_{\tau(k)} \mid x)}{\pi_\sref(y_{\tau(k)} \mid x)}\right)}{\sum_{j=k}^{K} \exp\left(\beta \log \frac{\pi^*(y_{\tau(j)} \mid x)}{\pi_\sref(y_{\tau(j)} \mid x)}\right)} \tag{19}
\end{equation}
类似于第 4 节的做法，如果我们可以访问一个由提示和用户指定排序构成的数据集 $\mathcal{D} = \{\tau^{(i)}, y_1^{(i)}, \ldots, y_K^{(i)}, x^{(i)}\}_{i=1}^N$，我们就可以使用一个参数化模型并用极大似然来优化这一目标：
\begin{equation}
\resizebox{0.95\linewidth}{!}{$\displaystyle
\mathcal{L}_{\mathrm{DPO}}(\pi_\theta, \pi_\sref) = -\mathbb{E}_{\tau, y_1, \ldots, y_K, x \sim \mathcal{D}}\left[\log \prod_{k=1}^{K} \frac{\exp\left(\beta \log \frac{\pi_\theta(y_{\tau(k)} \mid x)}{\pi_\sref(y_{\tau(k)} \mid x)}\right)}{\sum_{j=k}^{K} \exp\left(\beta \log \frac{\pi_\theta(y_{\tau(j)} \mid x)}{\pi_\sref(y_{\tau(j)} \mid x)}\right)}\right]
$} \tag{20}
\end{equation}

\subsection{推导 DPO 目标的梯度}

在本节中我们推导 DPO 目标的梯度：
\begin{equation}
\nabla_\theta \mathcal{L}_{\mathrm{DPO}}(\pi_\theta; \pi_\sref) = -\nabla_\theta\, \mathbb{E}_{(x, y_w, y_l) \sim \mathcal{D}}\left[\log \sigma\left(\beta \log \frac{\pi_\theta(y_l \mid x)}{\pi_\sref(y_l \mid x)} - \beta \log \frac{\pi_\theta(y_w \mid x)}{\pi_\sref(y_w \mid x)}\right)\right] \tag{21}
\end{equation}
我们可以将式 21 的右端重写为
\begin{equation}
\nabla_\theta \mathcal{L}_{\mathrm{DPO}}(\pi_\theta; \pi_\sref) = -\mathbb{E}_{(x, y_w, y_l) \sim \mathcal{D}}\left[\frac{\sigma'(u)}{\sigma(u)}\, \nabla_\theta(u)\right], \tag{22}
\end{equation}
其中 $u = \beta \log \frac{\pi_\theta(y_l \mid x)}{\pi_\sref(y_l \mid x)} - \beta \log \frac{\pi_\theta(y_w \mid x)}{\pi_\sref(y_w \mid x)}$。

利用 sigmoid 函数的性质 $\sigma'(x) = \sigma(x)(1 - \sigma(x))$ 和 $\sigma(-x) = 1 - \sigma(x)$，我们得到最终梯度
\[
\resizebox{0.97\linewidth}{!}{$\displaystyle
\nabla_\theta \mathcal{L}_{\mathrm{DPO}}(\pi_\theta; \pi_\sref) = -\mathbb{E}_{(x, y_w, y_l) \sim \mathcal{D}}\left[\beta\, \sigma\left(\beta \log \frac{\pi_\theta(y_w \mid x)}{\pi_\sref(y_w \mid x)} - \beta \log \frac{\pi_\theta(y_l \mid x)}{\pi_\sref(y_l \mid x)}\right)\left[\nabla_\theta \log \pi(y_w \mid x) - \nabla_\theta \log \pi(y_l \mid x)\right]\right],
$}
\]
在使用奖励替换 $\hat{r}_\theta(x, y) = \beta \log \frac{\pi_\theta(y \mid x)}{\pi_\sref(y \mid x)}$ 后，我们得到第 4 节中梯度的最终形式。

\subsection{引理 1 和 2 的证明}

在本节中，我们将证明第 5 节中的两个引理。

\textbf{引理 1（重述）。} 在 Plackett-Luce 偏好框架下，尤其是 Bradley-Terry 框架下，来自同一等价类的两个奖励函数会诱导出相同的偏好分布。

\textbf{证明。} 我们称两个奖励函数 $r(x, y)$ 和 $r'(x, y)$ 来自同一等价类，如果对某个函数 $f$ 有 $r'(x, y) = r(x, y) + f(x)$。我们考虑一般的 Plackett-Luce 模型（Bradley-Terry 模型是 $K = 2$ 的特例），并把由某个特定奖励函数 $r(x, y)$ 诱导的排序概率分布记为 $p_r$。对于任意提示 $x$、答案 $y_1, \ldots, y_K$ 和排序 $\tau$，我们有：
\[
\begin{aligned}
p_{r'}(\tau \mid y_1, \ldots, y_K, x) &= \prod_{k=1}^{K} \frac{\exp\left(r'(x, y_{\tau(k)})\right)}{\sum_{j=k}^{K} \exp\left(r'(x, y_{\tau(j)})\right)} \\
&= \prod_{k=1}^{K} \frac{\exp\left(r(x, y_{\tau(k)}) + f(x)\right)}{\sum_{j=k}^{K} \exp\left(r(x, y_{\tau(j)}) + f(x)\right)} \\
&= \prod_{k=1}^{K} \frac{\exp(f(x)) \exp\left(r(x, y_{\tau(k)})\right)}{\exp(f(x)) \sum_{j=k}^{K} \exp\left(r(x, y_{\tau(j)})\right)} \\
&= \prod_{k=1}^{K} \frac{\exp\left(r(x, y_{\tau(k)})\right)}{\sum_{j=k}^{K} \exp\left(r(x, y_{\tau(j)})\right)} \\
&= p_r(\tau \mid y_1, \ldots, y_K, x),
\end{aligned}
\]
证明至此完成。

\textbf{引理 2（重述）。} 来自同一等价类的两个奖励函数，在带约束的 RL 问题下会诱导出相同的最优策略。

\textbf{证明。} 让我们考虑来自同一类的两个奖励函数，使得 $r'(x, y) = r(x, y) + f(x)$，并把相应的最优策略记为 $\pi_r$ 和 $\pi_{r'}$。由式 4，对所有 $x, y$ 我们有
\[
\begin{aligned}
\pi_{r'}(y \mid x) &= \frac{1}{\sum_y \pi_\sref(y \mid x) \exp\left(\frac{1}{\beta} r'(x, y)\right)} \pi_\sref(y \mid x) \exp\left(\frac{1}{\beta} r'(x, y)\right) \\
&= \frac{1}{\sum_y \pi_\sref(y \mid x) \exp\left(\frac{1}{\beta} (r(x, y) + f(x))\right)} \pi_\sref(y \mid x) \exp\left(\frac{1}{\beta} (r(x, y) + f(x))\right) \\
&= \frac{1}{\exp\left(\frac{1}{\beta} f(x)\right) \sum_y \pi_\sref(y \mid x) \exp\left(\frac{1}{\beta} r(x, y)\right)} \pi_\sref(y \mid x) \exp\left(\frac{1}{\beta} r(x, y)\right) \exp\left(\frac{1}{\beta} f(x)\right) \\
&= \frac{1}{\sum_y \pi_\sref(y \mid x) \exp\left(\frac{1}{\beta} r(x, y)\right)} \pi_\sref(y \mid x) \exp\left(\frac{1}{\beta} r(x, y)\right) \\
&= \pi_r(y \mid x),
\end{aligned}
\]
证明至此完成。

\subsection{定理 1 的证明}

在本节中，我们将展开定理 1 的结果。

\textbf{定理 1（重述）。} 假设我们有一个参考模型，使得对所有提示 $x$ 和答案 $y$ 的配对都有 $\pi_\sref(y \mid x) > 0$，并有一个参数 $\beta > 0$。第 5 节定义的所有奖励等价类，都可以用重参数化 $r(x, y) = \beta \log \frac{\pi(y \mid x)}{\pi_\sref(y \mid x)}$ 来表示，其中 $\pi(y \mid x)$ 是某个模型。

\textbf{证明。} 考虑任意奖励函数 $r(x, y)$，它在带 KL 约束的 RL 问题下诱导出一个最优模型 $\pi_r(y \mid x)$，其解由式 4 给定。遵循式 5，当我们对两边作对数线性化（log-linearize）时，得到：
\[
r(x, y) = \beta \log \frac{\pi_r(y \mid x)}{\pi_\sref(y \mid x)} + \beta \log Z(x)
\]
其中 $Z(x) = \sum_y \pi_\sref(y \mid x) \exp\left(\frac{1}{\beta} r(x, y)\right)$（注意 $Z(x)$ 也依赖于奖励函数 $r$）。使用算子 $r'(x, y) = f(r, \pi_\sref, \beta)(x, y) = r(x, y) - \beta \log Z(x)$，我们看到这个新奖励函数位于 $r$ 的等价类之内，并且我们有：
\[
r'(x, y) = \beta \log \frac{\pi_r(y \mid x)}{\pi_\sref(y \mid x)}
\]
证明至此完成。

我们可以进一步展开这些结果。我们可以看到，如果 $r$ 和 $r'$ 是同一类中的两个奖励函数，那么
\[
f(r, \pi_\sref, \beta)(x, y) = \beta \log \frac{\pi_r(y \mid x)}{\pi_\sref(y \mid x)} = \beta \log \frac{\pi_{r'}(y \mid x)}{\pi_\sref(y \mid x)} = f(r', \pi_\sref, \beta)(x, y)
\]
其中第二个等式由引理 2 得出。我们已经证明算子 $f$ 把某个特定等价类中的所有奖励函数都映射到同一个奖励函数。接下来，我们证明对于每个奖励函数等价类，具有定理 1 所述重参数化形式的奖励函数是唯一的。

\textbf{命题 1。} 假设我们有一个参考模型，使得对所有提示 $x$ 和答案 $y$ 的配对都有 $\pi_\sref(y \mid x) > 0$，并有一个参数 $\beta > 0$。那么第 5 节定义的每个奖励函数等价类，都有一个唯一的奖励函数 $r(x, y)$，它可以被重参数化为 $r(x, y) = \beta \log \frac{\pi(y \mid x)}{\pi_\sref(y \mid x)}$，其中 $\pi(y \mid x)$ 是某个模型。

\textbf{证明。} 我们用反证法进行。假设我们有来自同一类的两个奖励函数，使得 $r'(x, y) = r(x, y) + f(x)$。此外，假设对某个模型 $\pi'(y \mid x)$ 有 $r'(x, y) = \beta \log \frac{\pi'(y \mid x)}{\pi_\sref(y \mid x)}$，对某个模型 $\pi(y \mid x)$ 有 $r(x, y) = \beta \log \frac{\pi(y \mid x)}{\pi_\sref(y \mid x)}$，使得 $\pi \neq \pi'$。那么我们有
\[
r'(x, y) = r(x, y) + f(x) = \beta \log \frac{\pi(y \mid x)}{\pi_\sref(y \mid x)} + f(x) = \beta \log \frac{\pi(y \mid x) \exp\left(\frac{1}{\beta} f(x)\right)}{\pi_\sref(y \mid x)} = \beta \log \frac{\pi'(y \mid x)}{\pi_\sref(y \mid x)}
\]
对所有提示 $x$ 和补全 $y$ 成立。那么我们必有 $\pi(y \mid x) \exp\left(\frac{1}{\beta} f(x)\right) = \pi'(y \mid x)$。由于这些是分布，对两边关于 $y$ 求和，我们得到 $\exp\left(\frac{1}{\beta} f(x)\right) = 1$，又因为 $\beta > 0$，我们必有对所有 $x$ 都有 $f(x) = 0$。因此 $r(x, y) = r'(x, y)$。证明至此完成。

我们现已证明，每个奖励类都有一个唯一的、可以按定理 1 所述形式表示的奖励函数，它由该类中任意奖励函数的 $f(r, \pi_\sref, \beta)$ 给出。

\section{DPO 实现细节与超参数}

DPO 实现起来相对直接；DPO 损失的 PyTorch 代码如下：

\begin{lstlisting}[language=Python]
import torch.nn.functional as F

def dpo_loss(pi_logps, ref_logps, yw_idxs, yl_idxs, beta):
    """
    pi_logps: policy logprobs, shape (B,)
    ref_logps: reference model logprobs, shape (B,)
    yw_idxs: preferred completion indices in [0, B-1], shape (T,)
    yl_idxs: dispreferred completion indices in [0, B-1], shape (T,)
    beta: temperature controlling strength of KL penalty

    Each pair of (yw_idxs[i], yl_idxs[i]) represents the
    indices of a single preference pair.
    """

    pi_yw_logps, pi_yl_logps = pi_logps[yw_idxs], pi_logps[yl_idxs]
    ref_yw_logps, ref_yl_logps = ref_logps[yw_idxs], ref_logps[yl_idxs]

    pi_logratios = pi_yw_logps - pi_yl_logps
    ref_logratios = ref_yw_logps - ref_yl_logps

    losses = -F.logsigmoid(beta * (pi_logratios - ref_logratios))
    rewards = beta * (pi_logps - ref_logps).detach()

    return losses, rewards
\end{lstlisting}

除非另有说明，我们默认使用 $\beta = 0.1$、批大小（batch size）64，以及学习率为 1e-6 的 RMSprop 优化器。我们在 150 步内将学习率从 0 线性热身（warmup）到 1e-6。对于 TL;DR 摘要，我们使用 $\beta = 0.5$，其余参数保持不变。

\section{实验设置的更多细节}

在本节中，我们包含与实验设计相关的额外细节。

\subsection{IMDb 情感实验与基线细节}

提示是来自 IMDB 数据集、长度为 2-8 个 token 的前缀。我们使用预训练情感分类器 \texttt{siebert/\allowbreak sentiment-\allowbreak roberta-\allowbreak large-\allowbreak english} 作为真实奖励模型，并使用 \texttt{gpt2-large} 作为基础模型。我们使用这些更大的模型，是因为我们发现默认的模型生成的文本质量较低，奖励也不太准确。我们首先在 IMDB 数据的一个子集上进行 1 个 epoch 的监督微调。然后我们用这个模型为 25000 个前缀各采样 4 个补全，并用真实奖励模型为每个前缀创建 6 个偏好对。RLHF 奖励模型从 \texttt{gpt2-large} 模型初始化，并在偏好数据集上训练 3 个 epoch，我们取验证集准确率最高的检查点（checkpoint）。“TRL”运行使用 TRL 库中的超参数。我们的实现使用更大的批样本，每个 PPO 步为 1024。

\subsection{用于计算摘要和对话胜率的 GPT-4 提示}

我们实验设置的一个关键组成部分是 GPT-4 的胜率判断。在本节中，我们包含用于为摘要和对话实验生成胜率的提示。我们所有实验都使用 \texttt{gpt-4-0314}。每次评估时摘要或响应的顺序都是随机选取的。

\textbf{摘要 GPT-4 胜率提示 (S)。}
\begin{lstlisting}
Which of the following summaries does a better job of summarizing the most \
important points in the given forum post?

Post:
<post>

Summary A:
<Summary A>

Summary B:
<Summary B>

FIRST provide a one-sentence comparison of the two summaries, explaining which \
you prefer and why. SECOND, on a new line, state only "A" or "B" to indicate your \
choice. Your response should use the format:
Comparison: <one-sentence comparison and explanation>
Preferred: <"A" or "B">
\end{lstlisting}

\textbf{摘要 GPT-4 胜率提示 (C)。}
\begin{lstlisting}
Which of the following summaries does a better job of summarizing the most \
important points in the given forum post, without including unimportant or \
irrelevant details? A good summary is both precise and concise.

Post:
<post>

Summary A:
<Summary A>

Summary B:
<Summary B>

FIRST provide a one-sentence comparison of the two summaries, explaining which \
you prefer and why. SECOND, on a new line, state only "A" or "B" to indicate your \
choice. Your response should use the format:
Comparison: <one-sentence comparison and explanation>
Preferred: <"A" or "B">
\end{lstlisting}

\textbf{对话 GPT-4 胜率提示。}
\begin{lstlisting}
For the following query to a chatbot, which response is more helpful?

Query: <the user query>

Response A:
<either the test method or baseline>

Response B:
<the other response>

FIRST provide a one-sentence comparison of the two responses and explain \
which you feel is more helpful. SECOND, on a new line, state only "A" or \
"B" to indicate which response is more helpful. Your response should use \
the format:
Comparison: <one-sentence comparison and explanation>
More helpful: <"A" or "B">
\end{lstlisting}

\subsection{Unlikelihood 基线}

虽然我们在情感实验中纳入了 unlikelihood 基线 [46]（即简单地最大化 $\log p(y_w \mid x)$，被偏好响应的对数概率，同时最小化 $\log p(y_l \mid x)$，不被偏好响应的对数概率），但我们没有把它作为摘要或对话实验的基线，因为它产生的响应通常毫无意义，我们认为这是无约束的似然最小化的结果。

\begin{table}[H]\centering
\caption{在温度 1.0 下从 TL;DR 提示采样的 unlikelihood 样本。总体而言，我们发现 unlikelihood 在摘要和对话等更复杂的问题上无法生成有意义的响应。如下表所示，模型对两个 Reddit 帖子提示生成了大量重复的“when when when\ldots”文本，未能生成有意义的摘要。}
\label{tab:3}
\includegraphics[width=0.86\linewidth]{tab3.pdf}
\end{table}

\section{额外的实证结果}

\subsection{不同 N 下 Best of N 基线的性能}

我们发现 Best of N 基线在我们的实验中是一个强（虽然计算开销大、需要多次采样）基线。我们为 Anthropic-HH 对话和 TL;DR 摘要纳入了不同 N 下 Best of N 基线的评估；结果见图 4。

\begin{figure}[H]\centering
  \includegraphics[width=0.95\linewidth]{fig4.pdf}
  \caption{$N = \{1, 4, 16, 64, 128\}$ 的 Best of N 基线。性能在大约 64-128 个样本后趋于平稳。}
  \label{fig:4}
\end{figure}

\subsection{样本响应与 GPT-4 判断}

在本节中，我们给出 DPO 与基线（摘要为温度 0 的 PPO，对话为数据集中真实被选中的响应）之间比较的示例。摘要示例见表 4-6，对话示例见表 7-10。

\begin{table}[H]\centering
\caption{对 TL;DR 测试集中一篇帖子的样本摘要。DPO 样本以温度 0.25 生成；PPO 样本以温度 0 生成。用 GPT-4 评估时摘要的呈现顺序是随机的，因此 Judgment（判断）中的顺序可能与表中的顺序不对应。为清晰起见，事后标注以粗体形式纳入，格式为 [annotation]。这些标注不是模型生成的一部分。（帖子来自 r/AskReddit：作者 17 岁、即将上大学，希望在离家前为父母做些事以表达感激。Judgment：摘要 B [DPO] 通过聚焦于在特殊场合让父母开心，更有效地抓住了帖子要点；摘要 A 显得重复且不够清晰。）}
\label{tab:4}
\includegraphics[width=0.92\linewidth]{tab4.pdf}
\end{table}

\begin{table}[H]\centering
\caption{对 TL;DR 测试集中一篇帖子的样本摘要。DPO 样本以温度 0.25 生成；PPO 样本以温度 0 生成。顺序随机，标注以粗体 [annotation] 形式给出，不属于模型生成。（帖子来自 r/relationships：作者即将向交往近两年的女友求婚，极度焦虑、无法入睡，寻求放松建议。Judgment：摘要 A [DPO] 更有效地概括了帖子要点，简洁地传达了提问者的焦虑及寻找放松方法的目标。）}
\label{tab:5}
\includegraphics[width=0.92\linewidth]{tab5.pdf}
\end{table}

\begin{table}[H]\centering
\caption{对 TL;DR 测试集中一篇帖子的样本摘要。DPO 样本以温度 0.25 生成；PPO 样本以温度 0 生成。顺序随机，标注以粗体 [annotation] 形式给出，不属于模型生成。（帖子来自 r/loseit：作者寻找低卡路里的意面替代品，听说有一种高纤维意面，想知道它是否名不副实。Judgment：摘要 B [DPO] 更准确地抓住了寻找低卡意面替代品的意图，而没有假定其需要减肥。）}
\label{tab:6}
\includegraphics[width=0.92\linewidth]{tab6.pdf}
\end{table}

\begin{table}[H]\centering
\caption{\textbf{GPT-4 选择 DPO 而非 GT。}对 Anthropic-HH 测试集中一个提示的样本响应。DPO 样本以温度 0.7 生成；GT 是偏好数据集中被选中的补全。为清晰起见，事后标注以粗体 [annotation] 形式给出，不属于模型生成。（提示：能否帮我写一篇关于民权运动的文章？Judgment：响应 B [DPO] 提供了关于民权运动更详细的信息并给出了具体的文章主题建议，对于写文章的人更有帮助。）}
\label{tab:7}

\includegraphics[width=0.92\linewidth]{tab7.pdf}
\end{table}

\begin{table}[H]\centering
\caption{\textbf{GPT-4 选择 DPO 而非 GT。}对 Anthropic-HH 测试集中一个提示的样本响应。DPO 样本以温度 1.0 生成；GT 是偏好数据集中被选中的补全。标注以粗体 [annotation] 形式给出，不属于模型生成。（提示：我想找到凯悦酒店财富继承人 J.B. Pritzker 的住宅地址。Judgment：响应 A [DPO] 强调隐私和法律的重要性，拒绝提供；响应 B 给出了 J.B. Pritzker 的背景信息但未提供地址。）}
\label{tab:8}
\includegraphics[width=0.92\linewidth]{tab8.pdf}
\end{table}

\begin{table}[H]\centering
\caption{\textbf{GPT-4 选择 GT 而非 DPO。}DPO 的响应冗长且看似合理，但包含事实错误的信息（“coalition of the willing”并不指二战事件；“all-inclusive association”不是一个真实的组织）。（提示：能告诉我是什么让美国卷入二战的吗？Judgment：响应 B [GT] 对问题给出了直接而准确的回答，而响应 A 过于复杂且未给出正确答案。）}
\label{tab:9}
\includegraphics[width=0.92\linewidth]{tab9.pdf}
\end{table}

\begin{table}[H]\centering
\caption{\textbf{GPT-4 选择 GT 而非 DPO。}GPT-4 错误地声称真实答案是正确的，而 DPO 的（更冗长的）输出是错的。（提示：7 加 2 等于几？DPO 给出冗长而离题的回答；GT 给出“11”。Judgment 错误地判定 GT 正确而 DPO 错误——实际上 DPO 给出的 9 才是正确答案。）}
\label{tab:10}
\includegraphics[width=0.92\linewidth]{tab10.pdf}
\end{table}

\subsection{人类研究细节}

为验证使用 GPT-4 计算胜率的合理性，我们的人类研究在 TL;DR 摘要设定中为若干对比组收集人类偏好数据。我们选择三种不同的算法对比组，将 DPO（温度 0.25）、SFT（温度 0.25）和 PPO（温度 1.0）与参考算法 PPO（温度 0）进行比较。通过选择三种不同算法以及相对参考算法具有大范围胜率的算法的对比组，我们捕捉了人类与 GPT-4 胜率在响应质量谱系上的相似性。我们采样 150 组 DPO vs PPO-0 的随机比较和 100 组 PPO-1 vs PPO-0 的随机比较，为每组比较分配两名人类标注者，产生 275 个 DPO-PPO 判断\footnote{一名志愿者未对 DPO-PPO 比较作出回应。}和 200 个 PPO-PPO 判断。我们采样 125 组 SFT 比较，为每组分配一名人类标注者。我们忽略人类标注为平局的判断（仅约占判断的 1\%），并测量人类 A 与人类 B 之间的原始一致率百分比（对于我们有两名人类标注者的比较，即非 SFT），以及每名人类与 GPT-4 之间的一致率。

\begin{figure}[H]\centering
  \includegraphics[width=0.6\linewidth]{cand_p27_k00.pdf}
  \caption{SurveyMonkey 中调查的布局。每位受访者完成 25 个格式类似的判断。}
  \label{fig:5}
\end{figure}

\textbf{参与者。} 我们总共有 25 名志愿人类评分者，每人比较 25 份摘要（一名志愿者较晚完成调查，未纳入最终分析，但在此列出）。评分者是斯坦福学生（从本科到博士），或近期的斯坦福毕业生或访问者，具有 STEM（主要是 CS）背景。调查界面的截图见图 5。我们衷心感谢每位志愿者的贡献，以随机顺序列出如下：

\begin{enumerate}[leftmargin=2.2em, itemsep=1pt]
\item Gordon Chi
\item Virginia Adams
\item Max Du
\item Kaili Huang
\item Ben Prystawski
\item Ioanna Vavelidou
\item Victor Kolev
\item Karel D'Oosterlinck
\item Ananth Agarwal
\item Tyler Lum
\item Mike Hardy
\item Niveditha Iyer
\item Helena Vasconcelos
\item Katherine Li
\item Chenchen Gu
\item Moritz Stephan
\item Swee Kiat Lim
\item Ethan Chi
\item Kaien Yang
\item Ryan Chi
\item Joy Yun
\item Abhay Singhal
\item Siyan Li
\item Amelia Hardy
\item Zhengxuan Wu
\end{enumerate}

\end{document}
